\chapter{牛顿运动定律}

在上一章中我们讨论了\uwave{质点运动学}，即如何描述质点的运动。

在这一章中我们将讨论\uwave{质点动力学}，即探究质点运动的原因。

\section{牛顿运动定律}

\subsection{牛顿第一定律}
\begin{BoxPrinciple}[牛顿第一定律]
    任何物体都保持静止或匀速直线运动状态，

    直到其他物体所作用于它的力迫使它改变这种状态为止。
\end{BoxPrinciple}
\Refx{定律}{牛顿第一定律}与两个基本的力学概念有关
\begin{enumerate}
    \item 我们将物体保持自身运动状态的原因，归结为物体的内部属性，称为\uwave{惯性}。
    \item 我们将物体运动状态发生改变的原因，归结为物体间相互作用，称为\uwave{力}。
\end{enumerate}

\Refx{定律}{牛顿第一定律}指出，物体自发具有保持运动状态不变的属性，物体的运动不需要力，力是改变运动状态的原因，惯性才是维持运动状态的原因，惯性是物体的内部属性。

应当指出，力是一个参照系无关的物理量，力描述物体间的相互作用，这种相互作用的大小和方向并不会随着观察所选取的参照系而变换，但运动学参量与参照系有关，相同的物体，在一个参照系中可能做匀速直线运动，在另一个参照系中可能就做匀变速直线运动。

这就导致牛顿第一定律并不总是正确的，其隐式的将参照系定义为两类，我们将牛顿第一定律成立的参照系称为\uwave{惯性系}，反之则称为\uwave{非惯性系}，参照系是否为惯性系需要通过实验来判断，即在参照系中通过力学实验，观察受力为零的物体是否满足加速度为零。

根据伽利略变换的结果
\begin{itemize}
    \item 一切相对于惯性系做匀速直线运动的参照系，均为惯性系。
    \item 一切相对于惯性系做变速直线运动的参照系，均为非惯性系。
\end{itemize}

实验表明，地面对于一般的力学实验是较为精确的惯性系，即在地面上进行的力学实验基本符合牛顿第一定律，但是由于地球存在自转，导致地面相对于地心具有很小的向心加速度，因此地面严格意义上是一个非惯性系，这可以通过傅科摆实验观测。

\subsection{牛顿第二定律}
牛顿第一定律阐释了力和运动的定性关系，牛顿第二定律将定量研究两者的关系，我们已经知道力的效果是使得物体的运动状态发生改变，那么，如果我们可以找到一个合适的物理量描述物体的“运动”\footnote{速度是不合适的，因为直观上感觉，速度为\SI{300}{m\cdot s^{-1}}的火箭似乎比相同速度的子弹具有强的多的“运动”。}，那么力就可以自然的定义为“运动的变化率”\footnote{这是合理的，例如加速度是“速度”改变的原因，而加速度也是“速度的变化率”，即关于时间的导数}。

牛顿在其巨著《自然哲学的数学原理》一书中，将牛顿第二定律表达为“运动变化”和力的关系，牛顿将他的叙述中的“运动”一词定义为物体质量和速度的乘积，这个物理量在现在被称为动量，记作$\vb*{P}=m\vb*{v}$，因此所谓的“运动变化”即指动量对时间的变化率。

牛顿第二定律使用现代语言可以表述为
\begin{BoxPrinciple}[牛顿第二定律]
    \centering\small
    物体的动量对时间的变化率与所加的合力成正比，

    且动量的方向与合力的方向相同。
\end{BoxPrinciple}
\Refx{定律}{牛顿第二定律}的数学形式为
\begin{BoxEquation}[牛顿第二定律的动量形式]
    \vb*{F}=\Fdif{\vb*{P}}{t}=\Fdif{\qty(m\vb*{v})}{t}
\end{BoxEquation}
力在国际单位制中的单位是\si{N=kg\cdot m\cdot s^{-2}}，称为\uwave{牛顿}，量纲是\si{[LMT^{-2}]}。

在经典力学中，物体的质量$m$被认为与时间无关，因此\Refx{公式}{牛顿第二定律的动量形式}中的质量可以作为常数提出，即$\vb*{F}=\Fdif{\vb*{P}}{t}=\Fdif{\qty(m\vb*{v})}{t}=m\Fdif{\vb*{v}}{t}=m\vb*{a}$，即
\begin{BoxEquation}[牛顿第二定律的加速度形式]
    \vb*{F}=m\vb*{a}
\end{BoxEquation}
应当指出，尽管\Refx{公式}{牛顿第二定律的加速度形式}是更为被熟知的形式，但是\Refx{公式}{牛顿第二定律的动量形式}事实才是更为一般的形式，例如研究物体质量随时间变化的物体的运动\footnote{类似于火箭在运动中随着燃料消耗质量逐渐减小的过程，并非指由相对论因素导致的质量变化，}，例如在相对论或量子力学中动量不再由质量和速度的乘积定义时，动量形式的牛顿第二定律仍然是可以适用的，但此时$\vb*{F}=m\vb*{a}$往往就不再成立了。

根据\Refx{公式}{牛顿第二定律的加速度形式}，如果相同的力作用在两个质量分别为$m_1,m_2$的物体上，设产生的加速度的大小分别为$a_1,a_2$，那么有
\begin{equation*}
    \frac{m_1}{m_2}=\frac{a_2}{a_1}
\end{equation*}
即在相同的力的作用下，质量和加速度成反比，质量大的物体加速度小，这就意味着质量大的物体抵抗运动变化能力强，因此此处的质量是物体惯性的量度，称为\uwave{惯性质量}。

实验表明，当作用于质点上的力有多个时，其合力可以表达为分力的矢量和
\begin{BoxEquation}[力的叠加原理]
    \vb*{F}=\vb*{F_1}+\vb*{F_2}+\cdots+\vb*{F_n}=\Sum{i=1}{N}{\vb*{F_i}}
\end{BoxEquation}

\Refx{定律}{牛顿第二定律}仅适用于惯性系，且在所有惯性系中定律的形式不变（因此我们不可能通过力学实验确定惯性系的运动速度），换言之，惯性系是平等的，不存在一个特殊的惯性系，因此绝对静止的参考系的不存在的，这个结论称为\uwave{力学相对性原理}。

\Refx{定律}{牛顿第二定律}是经典力学的核心，但注意其仅适用于宏观低速的运动。

\subsection{牛顿第三定律}
\begin{BoxPrinciple}[牛顿第三定律]
    \centering\small
    两个物体之间的相互作用力大小相等，方向相反且沿同一直线

    分别作用在两个物体上。
\end{BoxPrinciple}
\Refx{定律}{牛顿第三定律}说明力作为物体间的相互作用，应具有的基本性质。
\begin{enumerate}
    \item 成对性：作用力和反作用力总是成对出现，因为物体间的作用是相互的。
    \item 同时性：作用力和反作用力总是同时产生，同时消失，没有因果之分。
    \item 一致性：作用力和反作用力是性质相同的力。
\end{enumerate}
\Refx{定律}{牛顿第三定律}中，需要辨析“作用力与反作用力”和“平衡力”的概念，例如置于桌子的木块受到两个力，地球对木块的引力，桌子对木块的弹力，两者平衡，是一对平衡力，但前者的反作用力是木块对地球的引力，而后者的反作用力是木块对桌子的弹力。

\Refx{定律}{牛顿第三定律}的意义在于其揭示了力的一个重要特征——力具有实质性起源的，这就使得力不仅仅是一个“动量变化率”的定义，当观察到一个物体的动量发生变化时，我们在断言这是力的作用之外，我们还一定能在物体附近找到某个物体是力的来源。

\Refx{定律}{牛顿第三定律}在惯性系和非惯性系中均是成立的，因为力是参照系无关的，而第三定律只涉及力自身的性质，不同于第一第二定律包含力和运动的关系，因此普遍成立。

\Refx{定律}{牛顿第三定律}只保证对于彼此接触的物体才适用，对于超距作用的力，其不一定成立，因为我们知道，诸如引力和电磁力这样的超距作用，本质是通过场的方式进行的，例如日地间的引力，实际是由“地球与引力场的作用”以及“太阳与引力场的作用”两部分组成，并不是牛顿第三定律条件中所要求的“物体间直接的相互作用”。

不过幸运的是，对于不需要考虑相对论效应的引力，以及静止电荷间的电磁力，我们可以近似的认为其也是符合牛顿第三定律的（运动电荷间的电磁力就不完全符合）。

\section{力学中的常见力}
在开始这一节的内容前，我们首先强调，尽管牛顿定律是经典力学的基石，但牛顿定律自身其实并没有揭示出力的全部性质，或者说孤立状态的牛顿定律是无意义的。
\begin{itemize}
    \item 我们尝试判断一个参照系是否为惯性系，那我们需要通过实验观察静止的物体是否受力均为零，但是什么是受力为零，这需要通过牛顿第二定律给出，但是牛顿第二定律可以使用的前提是参照系是一个惯性系，这就构成了循环依赖
    \begin{center}
        \small
        判断是否为惯性系需要牛顿第二定律\quad$\Longleftrightarrow$\quad 只有在惯性系中牛顿第二定律才是成立的
    \end{center}
    这就意味着不可能仅通过牛顿定律，判断参照系是否为惯性系。
    \item 我们即便可以通过某种方式找到一个已知的惯性系，但此时的力也是毫无意义的
    \begin{itemize}
        \item 我们宣称力是加速度和质量的乘积，并正式的加以定义，因此我们可以得出结论，如果我们已知物体的受力，那么物体的加速度就是力除以物体的质量。
        \item 我们将运动状态改变的原因归结为力，那么如果我们观察到某个物体保持静止或做匀速直线运动，那我们可以断言，这个物体一定是不受力的。
    \end{itemize}
    很容易看出上面的内容是一些毫无意义的循环论证，我们从中没有任何发现，正如费恩曼所言，我们可以整天坐在安乐椅上，任意定义一些词，例如我们可以说
    \begin{center}
        \small
        
        物体总是具有自发保持静止的属性。
        
        物体运动的原因可以归结于戈斯\footnote{戈斯即Gorce，是力Force的改写的音译。}的作用，戈斯是位置的变化率\footnote{戈斯就是物体的速度，速度为零物体静止，速度非零物体运动，戈斯的存在似乎毫无意义，这也正是这个例子想表达的含义。}。
    \end{center}
    从中我们可以得到一条十分伟大的定律：一切物体都将保持静止，直到有戈斯作用在物体上。这就与牛顿第一定律在形式上非常相似，但很明显，这并不包含任何内容。
\end{itemize}
以上讨论就意味着，牛顿定律是不完全的定律，力必然具有某种独立于牛顿定律的特有性质，牛顿定律暗示，如果我们研究质量和加速度的乘积构成的物理量，将该物理量称为力，进而研究力的特征，我们会发现力具有某种简单性，这对于定量分析是极有帮助的。

下面我们将揭示力在牛顿定律之外遵循的规律，力的类型是多样的，力遵循的规律也是多样的，在这里我们主要介绍力学中三种最常见的力：万有引力，弹性力，摩擦力。

\subsection{万有引力}
\uwave{万有引力}是存在于任何两个物体之间的吸引力，是自然界中的基本力之一。

\begin{BoxPrinciple}[万有引力定律]
    任何两个质点之间都存在相互作用的引力，力的方向沿着两质点的连线，
    
    力的大小与两质点的质量的乘积成正比，与两质点之间的距离的平方成反比。
    \begin{BoxEq}
        F=G\frac{m_1m_2}{r^2}
    \end{BoxEq}
\end{BoxPrinciple}
其中\uwave{引力常量}$G=6.67\times 10\times 10^{-11}$\si{N\cdot m^2\cdot kg^{-2}}。

\Refx{定律}{万有引力定律}中引入的质量称为\uwave{引力质量}，引力质量是物体之间相互吸引强弱的量度，\Refx{公式}{牛顿第二定律的加速度形式}中的质量称为\uwave{惯性质量}，惯性质量是改变物体运动状态的难易程度的度量，两者是不同的，两种质量分别代表物体的两种不同属性，通俗的说，对于一个物体，\textbf{它称起来有多重是一回事，它运动起来有多难是另一回事}。

但精确的实验研究和理论分析表明，对于任何物体来说，其引力质量和惯性质量是相等的，因此我们可以不加区分的使用引力质量和惯性质量，即将两者当成完全相同的概念“质量”。

如果在\Refx{定律}{万有引力定律}中，令$m_1=m_e$代表地球的质量，令$m_2=m$代表地表上某一物体的质量，令$r=R_e$代表地球的半径，那么联立牛顿第二定律可得
\begin{equation*}
    F=G\frac{m_em}{R_e^2}=ma
\end{equation*}
上式中，左侧的质量$m$为地表物体的引力质量，右式的质量$m$为地表物体的惯性质量，先前我们讨论过，这两种质量可以不加区分，因此我们可以将两者约去，并改记$a$为$g$
\begin{BoxEquation}[重力加速度]
    g=G\frac{m_e}{R_e^2}
\end{BoxEquation}
上式中的$g$称为\uwave{重力加速度}，其方向竖直向下指向地心，由于地表附近的物体到地心的距离均可以大致近似为地球半径，因此重力加速度可以视为一个常数，通常取$g=9.8$\si{m\cdot s^{-2}}。
\begin{itemize}
    \item 重力加速度的数值与物体本身的质量无关。
    \item 重力加速度会随着物体高度增加而减小，但这种变化在地表附近可以被忽略。
\end{itemize}
由于地表附近的重力加速度为常数，那么此时的引力可以近似为$mg$，称为\uwave{重力}，记为
\begin{BoxEquation}[重力公式]
    P=-mg
\end{BoxEquation}
上式的负号代表重力方向总是竖直向下的。

\subsection{弹性力}
\uwave{弹性力}也称为\uwave{弹力}，是物体受外力作用形变后产生的一种回复力，是一种接触力。

弹性力主要有以下三种表现形式
\begin{enumerate}
    \item 第一种是两物体通过一定的面积相互接触的情况，这时相互压紧的物体都会发生微小形变，因而产生对对方的弹力作用，这种弹力通常称为\uwave{支持力}。
    
    支持力的大小取决于压紧程度，支持力总是垂直于接触面指向对方。

    \item 第二种是绳或线对物体的牵拉作用，这种作用是由于绳发生了微小的形变而产生的，在外部，绳将对物体产生\uwave{拉力}，在内部，绳的各段之间会产生相互作用的\uwave{张力}。
    
    拉力和张力的大小取决于绳的拉紧程度，拉力（绳对物体的）总是指向绳收缩的方向。

    \item 第三种是弹簧的弹性力，不同于支持力和拉力的大小是完全取决于外部条件，弹簧的弹性力只取决于弹簧的性质和形变状况，弹簧的弹性力在弹性限度内满足\uwave{胡克定律}
\end{enumerate}
\begin{BoxPrinciple}[胡克定律]
    \begin{BoxEq}
        F=-kx
    \end{BoxEq}
\end{BoxPrinciple}
上式的负号代表弹簧弹力的方向总是与形变方向相反，指向平衡位置。

上式的$k$称为弹簧的\uwave{颈度系数}，这取决于弹簧的结果，而$x$代表弹簧形变的位移。

\subsection{摩擦力}
\uwave{摩擦力}是指当两个相互接触的物体，发生相对运动，或具有相对运动趋势时，这两个物体的接触面上产生的平行于接触面，阻碍物体相互运动的相互作用力。

当相互接触的两物体，在外力作用下，并没有产生相对运动，此时我们就将阻碍相对运动趋势的原因归结于\uwave{静摩擦力}，静摩擦力的大小视外力的大小而定，总是使物体保持相对静止，静摩擦力具有一个最大限度值，称为\uwave{最大静摩擦力}，实验表明其满足公式
\begin{BoxEquation}[最大静摩擦力公式]
    F_s=\mu_sN
\end{BoxEquation}
其中$N$代表接触面上的弹力，而$\mu_s$称为\uwave{静摩擦因数}，通常与接触面的材料有关。

当相互接触的两物体所受到的外力超过最大静摩擦力时，物体间就会产生相对运动，但此时仍有阻碍相对运动的力，只不过这种阻碍小于外力，称为\uwave{滑动摩擦力}，满足公式
\begin{BoxEquation}[滑动摩擦力公式]
    F_f=\mu_fN
\end{BoxEquation}
其中$N$代表接触面上的弹力，而$\mu_f$称为\uwave{动摩擦因数}，除了与接触面的材料有关，通常还与相对速度有关，在大多数情况下随着速度增大而减小，此外对于一定的接触面有$\mu_s>\mu_f$，也就是说在摩擦力的阻碍下，使物体开始运动需要耗费比维持物体运动更大的力。

通常情况下，静摩擦因数$\mu_s$和动摩擦因数$\mu_f$的值均小于$1$。

在简要的分析中，我们可以认为动摩擦因数$\mu_f$是与速率无关的常数，同时可以假定动摩擦因数与静摩擦因数的值相同，即有$\mu=\mu_f=\mu_s$（简称为\uwave{摩擦因数}）。

摩擦力和弹性力与先前的万有引力是不同的。万有引力是一种基本作用力，这意味随着对引力的研究深入，其数学形式将变得愈发简洁，例如“使苹果落地的引力”和“使月球绕地球旋转的引力”均可以统一在相同的万有引力公式中，两者均遵循平方反比规律。

摩擦力和弹力则并不会如此，考虑的越细致，我们会发现其规律愈发复杂，例如，简单的分析中摩擦力仅与弹力大小有关，复杂的分析下我们还要考虑其和速度的关系，等等。

这是因为我们对摩擦力和弹力的大小刻画，是一种基于工程的经验公式。两者的本质是一种称为电磁力的基本作用力，后者与万有引力一样具有简单性，但在实际问题我们难以根据电磁力的规律推出摩擦力和弹力的大小，因此这些经验公式是非常有必要的。

\subsection{非惯性系下的牛顿第二定律}
我们知道，牛顿第二定律是普遍适用于惯性系，对于两个惯性系$K_1,K_2$，根据\Refx{公式}{伽利略加速度变换}可知$\vb*{a}_2=\vb*{a}_1$，再运用\Refx{公式}{牛顿第二定律的加速度形式}
\begin{equation*}
    \vb*{F}=m\vb*{a}_2=m\vb*{a}_1
\end{equation*}
关注到在两个坐标系中根据牛顿第二定律计算得到的力是相同的，这符合力的大小与参照系无关的假设，因此，牛顿定律对于一切惯性系均具有相同的形式，这就意味着，力学规律具有\uwave{空间平移对称性}，如果再考虑\Refx{公式}{伽利略时间变换}，其还具有\uwave{时间平移对称性}。

然而，如果$K_1$是惯性系，而$K_2$是非惯性系，那么由式\ref{式:伽利略加速度变换的扩充}就有$$\vb*{a}_2=\vb*{a}_1-\vb*{a}_x$$
其中$\vb*{a}_x$为非惯性系$K_2$相对于惯性系$K_1$的加速度。

在$K_1$中应用牛顿第二定律
\begin{equation*}
    \vb*{F}=m\vb*{a}_1=m\vb*{a}_2+m\vb*{a}
\end{equation*}
在$K_2$中应用牛顿第二定律
\begin{equation*}
    \vb*{F}=m\vb*{a}_2\quad\text{（错误！）}
\end{equation*}
关注到上述两个式子发生了矛盾，这是因为牛顿第二定律对于非惯性系$K_2$是不成立的，这就是说，在非惯性系$K_2$中的加速度$a_2$与力的关系不遵循$\vb*{F}=ma_2$，根据惯性系$K_1$中得出的结论，我们知道两者遵循的正确规律实际应为$\vb*{F}-m\vb*{a}_x=m\vb*{a}_2$。

那么为了简洁表达非惯性系中的动力学规律，我们引入\uwave{惯性力}
\begin{BoxEquation}[非惯性系中的惯性力]
    \vb*{F}_x=-m\vb*{a}_x
\end{BoxEquation}
那么非惯性系中的动力学规律就可以表达为（记非惯性系中的加速度$\vb*{a}_2=\vb*{a}$）
\begin{BoxEquation}[非惯性系中的动力学规律]
    \vb*{F}+\vb*{F}_x=m\vb*{a}
\end{BoxEquation}
这就得到了一个类似于牛顿第二定律的结果，右式仍然为该参照系的加速度和质量的乘积，左式则在合外力的基础上添加了一个惯性力，但应当注意，惯性力是虚拟的力，\textbf{惯性力的本质是将非惯性系与惯性系的差异比拟为力的作用，从而模拟相对运动的动力学结果}，因此它不是物体间的相互作用，没有施力物体，没有反作用力，不满足牛顿第三定律。
